\documentclass{book}

\title{Physical Computing}
\newcommand{\booksubtitle}{The Principles of Computing for Physicists}
\newcommand{\booklicense}{Creative Commons Zero 1.0 Universal}

\author{Ed Daw}
% Author subtitle could be a university or a geographical location, for example
\newcommand{\authorsubtitle}{PHY31001, The University of Sheffield}

\newcommand{\VHDL}{\texttt{VHDL}\;}
\newcommand{\git}{\texttt{git}\;}
\newcommand{\SVN}{\texttt{SVN}\;}
\newcommand{\CVS}{\texttt{CVS}\;}
\newcommand{\github}{\texttt{GitHub}\;}
\newcommand{\vhdl}{\texttt{VHDL}\;}

% Create convenient commands \booktitle and \bookauthor
\makeatletter
\newcommand{\booktitle}{\@title}
\newcommand{\bookauthor}{\@author}
\makeatother

% This utf8 declaration is not needed for versions of latex > 2018 but may
% be helpful for older software. Eventually it may not be worth keeping.
\usepackage[utf8]{inputenc}  
\usepackage{fix-cm}  % this package allows large \fontsize
\usepackage{tikz}    % this is for graphics. e.g. rectangle on title page
\usepackage{amsmath} % Used by equations
\usepackage{halloweenmath}
% used for quantum mechanics
\usepackage{braket}
\usepackage{minted}
\usepackage{appendix}
\usepackage{xurl}
\usepackage{hyperref}
\usepackage{rotating}

% The following dimensions specify 4.75" X 7.5" content on 6 3/8" by 9 1/4"
% paper. The paper width and height can be tweaked as required and the content
% should size to fit within the margins accordingly.
%
% The (inside) bindingoffset should be larger for books with more pages. Some
% standard recommended sizes are .375in minimum up to 1in for 600+ page books.
% Sizes .75in and .875in are also recommended roughly at 150 and 400 pages.
\usepackage[bindingoffset=0.625in,
            left=.5in, right=.5in,
            top=.8125in, bottom=.9375in,
            paperwidth=6.375in, paperheight=9.25in]{geometry}
% Here is an alternative geometry for reading on letter size paper:
% \usepackage[margin=.75in, paperwidth=8.5in, paperheight=11in]{geometry}

\renewcommand{\contentsname}{Table of Contents} % default is {Contents}
\usepackage{makeidx}

% Divide book soource into subfiles
\usepackage{subfiles}

\makeindex % Initialize an index so we can add entries with \index

% Content Starts Here
\begin{document}
\frontmatter

% ---- Half Title Page ----
% current geometry will be restored after title page
\newgeometry{top=1.75in,bottom=.5in}
\begin{titlepage}
\begin{flushleft}

% Title
\textbf{\fontfamily{qcs}\fontsize{48}{54}\selectfont Physical\\Computing\\}

% Draw a line 4pt high
\par\noindent\rule{\textwidth}{4pt}\\
Sequential
% Subtitle
% Shaded box from left to right. Text node is midway (centered).
\begin{tikzpicture}
\shade[bottom color=lightgray,top color=white]
    (0,0) rectangle (\textwidth, 1.5)
    node[midway] {\textbf{\large \textit{\booksubtitle}}};
\end{tikzpicture}

% Edition Number
\begin{flushright}
\Large First Edition
\end{flushright}

% \vspace{\fill}
\vspace{\fill}

\end{flushleft}
\end{titlepage}
\restoregeometry
% ---- End of Half Title Page ----


% No page numbers on the Frontispiece page
\thispagestyle{empty}


% ---- Title Page ----
% current geometry will be restored after title page
\newgeometry{top=1.75in,bottom=.5in}
\begin{titlepage}
\begin{flushleft}

% Title
\textbf{\fontfamily{qcs}\fontsize{48}{54}\selectfont Physical\\Computing\\}

% Draw a line 4pt high
\par\noindent\rule{\textwidth}{4pt}\\

% Shaded box from left to right with Subtitle
% The text node is midway (centered).
\begin{tikzpicture}
\shade[bottom color=lightgray,top color=white]
    (0,0) rectangle (\textwidth, 1.5)
    node[midway] {\textbf{\large \textit{\booksubtitle}}};
\end{tikzpicture}

% Edition Number
\begin{flushright}
\Large First Edition
\end{flushright}

\vspace{\fill}

% Author and Location
\textbf{\large \bookauthor}\\[3.5pt]
\textbf{\large \textit{\authorsubtitle}}

\vspace{\fill}

% Self Publishing Logo. Free to use: CC0 license.
% The source file is book.svg. If you change the svg, you must then convert
% it to pdf. There are many online and offline tools available to do that.
\begin{center}
\includegraphics{booksvg.pdf}\\[4pt]
\fontfamily{lmtt}\small{Self Publishers Worldwide\\
Seattle San Francisco New York\\
London Paris Rome Beijing Barcelona}
\end{center}

\end{flushleft}
\end{titlepage}
\restoregeometry
% ---- End of Title Page ----

% Do not show page numbers on colophon page
\thispagestyle{empty}

\begin{flushleft}
\vspace*{\fill}
This book was typeset using \LaTeX{} software.\\
\vspace{\fill}
Copyright \textcopyright{} \the\year{}  \bookauthor\\
License: \booklicense
\end{flushleft}

% A title page resets the page # to 1, but the second title page
% was actually page 3. So add two to page counter.
\addtocounter{page}{2}

\subfile{preface/physical_computing_preface.tex}

% Three-level Table of Contents
\setcounter{tocdepth}{3}
\tableofcontents

\mainmatter

\subfile{chapter_1/physical_computing_chapter_1.tex}

\subfile{chapter_2/physical_computing_chapter_2.tex}

\subfile{chapter_3/physical_computing_chapter_3.tex}

\subfile{chapter_4/physical_computing_chapter_4.tex}

%\chapter{Serial Buses }
\chapter{Data Converters}
\chapter{State Machines}
\chapter{Instruction Sets and Computers}
\chapter{Programming in C}
\chapter{Buses again}
\chapter{A Command Line Interpreter}
\chapter{Qubits and Quantum Computing}

\appendix
\appendixpage
\noappendicestocpagenum
\addappheadtotoc

%\subfile{appendix_3/physical_computing_appendix_3.tex}
%\subfile{appendix_4/physical_computing_appendix_4.tex}
\subfile{appendix_5/physical_computing_appendix_5.tex}
\subfile{appendix_6/physical_computing_appendix_6.tex}
\subfile{appendix_7/physical_computing_appendix_7.tex}
\subfile{appendix_8/physical_computing_appendix_8.tex}
\subfile{appendix_9/physical_computing_appendix_9.tex}
\subfile{appendix_10/physical_computing_appendix_10.tex}

\bibliographystyle{plain}
\bibliography{physical_computing}

\backmatter
\addcontentsline{toc}{chapter}{Index}
\printindex
\end{document}